\documentclass[11pt,a4paper]{article}

\usepackage[utf8]{inputenc}
\usepackage[T1]{fontenc}
\usepackage[brazil]{babel}
\usepackage[a4paper,margin=2.5cm]{geometry}
\usepackage{lmodern}
\usepackage{microtype}
\usepackage{setspace}
\usepackage{graphicx}
\usepackage{xcolor}
\usepackage{booktabs}
\usepackage{array}
\usepackage{enumitem}
\usepackage{tikz}
\usetikzlibrary{arrows.meta,positioning,shapes.geometric,fit,backgrounds}
\usepackage[hidelinks]{hyperref}

\onehalfspacing
\setlength{\parindent}{1.25cm}
\setlength{\parskip}{0.15em}

\definecolor{azul}{RGB}{26,76,122}
\definecolor{verde}{RGB}{42,122,77}
\definecolor{cinza}{RGB}{242,245,247}

\title{\textbf{Estacionamentos inteligentes utilizando uma Rede de Sensores Sem Fio}}
\author{Álvaro Lucio Almeida Ribeiro\\Matheus Henrique Fonseca Afonso}
\date{2026}

\begin{document}

\begin{titlepage}
\centering
\vspace*{2cm}
{\large \textbf{TRABALHO DE PESQUISA}}\\[2cm]
{\LARGE \textbf{Estacionamentos inteligentes utilizando uma Rede de Sensores Sem Fio}}\\[1.8cm]
{\large Álvaro Lucio Almeida Ribeiro\\[0.2cm]Matheus Henrique Fonseca Afonso}\\[0.8cm]
\vfill
\begin{minipage}{0.82\textwidth}
\centering
Este relatório apresenta uma proposta simples de Rede de Sensores Sem Fio (RSSF/WSN) para identificar vagas livres em estacionamentos e disponibilizar essa informação em um aplicativo.
\end{minipage}
\vfill
{\large 2026}
\end{titlepage}

\pagenumbering{arabic}
\section*{Resumo}
\addcontentsline{toc}{section}{Resumo}
Encontrar uma vaga em regiões movimentadas pode aumentar o tempo de deslocamento, o consumo de combustível e a emissão de poluentes. Este trabalho pesquisa a aplicação de uma Rede de Sensores Sem Fio (RSSF) em estacionamentos inteligentes. A proposta utiliza um nó sensor em cada vaga para detectar a presença de um veículo, uma rede de comunicação de baixo consumo para transportar as leituras até um gateway e um serviço de software que atualiza um aplicativo. São discutidas alternativas de sensores, tecnologias de comunicação, organização da rede, consumo de energia, confiabilidade e segurança. A RSSF é útil porque permite instalar muitos pontos de medição distribuídos sem exigir cabeamento individual entre todas as vagas, além de possibilitar expansão e manutenção por partes.

\noindent\textbf{Palavras-chave:} rede de sensores sem fio; estacionamento inteligente; Internet das Coisas; LoRa; Zigbee.

\section{Introdução}
O estacionamento é um problema cotidiano em centros urbanos, universidades, hospitais e shopping centers. O motorista normalmente percorre os corredores sem saber quais vagas estão disponíveis. Quando muitos veículos fazem isso ao mesmo tempo, ocorre um fluxo desnecessário de circulação: aumenta o tempo de procura, formam-se filas e são produzidas emissões que poderiam ser evitadas.

Um estacionamento inteligente transforma a ocupação das vagas em informação digital. Sensores instalados nas vagas identificam se há ou não um veículo; a rede transmite as leituras; e um sistema apresenta no aplicativo um mapa ou uma lista de vagas livres. O usuário pode então escolher uma vaga antes de entrar ou durante o deslocamento dentro do estacionamento.

A aplicação escolhida neste trabalho é, portanto, o monitoramento de vagas de estacionamento com RSSF. A escolha é adequada porque o cenário possui muitos pontos de observação, distribuídos espacialmente, e cada ponto transmite uma pequena quantidade de dados. Esses são requisitos típicos de uma rede de sensores: nós compactos, baixo consumo, comunicação sem fio e processamento distribuído. Trabalhos anteriores descrevem arquiteturas com sensores nas vagas, nós repetidores e um gateway conectado à Internet \cite{bagula2015,chen2016}.

\section{Por que utilizar uma RSSF?}
Uma RSSF é formada por vários nós sensores capazes de medir uma grandeza física, processar uma leitura e enviá-la por rádio. No estacionamento, a grandeza de interesse é a ocupação da vaga. O dado pode ser representado de maneira simples como $ocupada = 0$ ou $ocupada = 1$, acompanhado do identificador da vaga e do horário da medição.

Há cinco vantagens principais para esse cenário:
\begin{enumerate}[label=\arabic*.,leftmargin=1cm]
    \item \textbf{Instalação flexível:} os sensores podem ser distribuídos pelo piso, teto ou laterais, sem levar cabos de dados até cada vaga.
    \item \textbf{Escalabilidade:} novas vagas podem ser incluídas adicionando novos nós e cadastrando seus identificadores no sistema.
    \item \textbf{Baixo volume de dados:} uma vaga envia apenas eventos de mudança ou mensagens periódicas de estado, economizando energia e largura de banda.
    \item \textbf{Manutenção localizada:} uma falha geralmente afeta um nó ou um pequeno trecho, sem interromper necessariamente todo o estacionamento.
    \item \textbf{Integração com serviços:} o gateway pode enviar as leituras para uma API, um banco de dados e um aplicativo móvel.
\end{enumerate}

Além de ajudar o motorista, a solução fornece ao administrador indicadores de ocupação, horários de pico e tempo médio de permanência. Esses dados podem apoiar a gestão de vagas reservadas, acessíveis ou destinadas a veículos elétricos. A literatura também destaca que a posição dos nós influencia a cobertura e o tempo de vida da rede, tornando o planejamento parte importante do projeto \cite{bagula2015}.

\section{Arquitetura proposta}
A arquitetura sugerida possui quatro camadas: sensoriamento, comunicação, serviço e aplicação. Cada vaga recebe um nó sensor, composto por sensor de presença, microcontrolador, rádio e fonte de energia. Os nós enviam dados diretamente ao gateway ou por saltos entre nós vizinhos. O gateway concentra as mensagens e usa Wi-Fi, Ethernet ou rede celular para alcançar o servidor.

\begin{figure}[ht]
\centering
\begin{tikzpicture}[
node distance=0.65cm and 0.85cm,
box/.style={draw, rounded corners, align=center, minimum height=1.05cm, minimum width=2.55cm, fill=cinza},
sensor/.style={box, fill=blue!8, draw=azul},
gateway/.style={box, fill=orange!15, draw=orange!70!black},
service/.style={box, fill=green!10, draw=verde},
app/.style={box, fill=purple!8, draw=purple!60!black},
arrow/.style={-{Latex[length=2mm]}, thick, draw=azul}
]
\node[sensor] (s1) {Nó da vaga 1\\sensor + rádio};
\node[sensor, right=of s1] (s2) {Nó da vaga 2\\sensor + rádio};
\node[sensor, right=of s2] (s3) {Nó da vaga $n$\\sensor + rádio};
\node[gateway, below=1.1cm of s2] (gw) {Gateway\\agregação};
\node[service, below=1.1cm of gw] (sv) {Servidor/API\\banco de dados};
\node[app, below=1.1cm of sv] (ap) {Aplicativo\\mapa de vagas};
\draw[arrow] (s1) -- (gw);
\draw[arrow] (s2) -- (gw);
\draw[arrow] (s3) -- (gw);
\draw[arrow] (gw) -- node[right,font=\scriptsize]{Internet} (sv);
\draw[arrow] (sv) -- (ap);
\end{tikzpicture}
\caption{Arquitetura conceitual da RSSF para estacionamento inteligente.}
\label{fig:arquitetura}
\end{figure}

\subsection{Funcionamento de um nó sensor}
O microcontrolador mantém o nó em modo de baixo consumo e acorda em intervalos definidos. Ao acordar, realiza a leitura, aplica uma filtragem simples e compara o resultado com um limiar. Quando o estado muda de livre para ocupado, ou de ocupado para livre, o nó transmite um evento. Uma mensagem pode ter o formato: \texttt{ID\_VAGA, ESTADO, BATERIA, HORARIO}.

Uma confirmação do gateway evita que uma perda de pacote deixe o aplicativo com informação antiga. Se não houver confirmação, o nó tenta reenviar a mensagem com um intervalo limitado. Mensagens periódicas de ``estou funcionando'' (heartbeat) permitem identificar sensores sem comunicação. Para evitar falsos eventos, é recomendável exigir duas ou mais leituras consecutivas com o mesmo resultado antes de alterar o estado da vaga.

\subsection{Comunicação e topologia}
Em um estacionamento pequeno, uma topologia estrela pode ser suficiente: todos os nós comunicam diretamente com o gateway. Em estacionamentos grandes ou subterrâneos, uma topologia em árvore ou malha pode usar nós repetidores para contornar obstáculos. A seleção da tecnologia depende da distância, do número de vagas, da existência de energia elétrica e da necessidade de latência.

\begin{table}[ht]
\centering
\caption{Alternativas de comunicação para a RSSF.}
\label{tab:comunicacao}
\begin{tabular}{p{2.2cm}p{3.7cm}p{4.5cm}}
\toprule
Tecnologia & Característica principal & Uso sugerido \\
\midrule
Zigbee/IEEE 802.15.4 & Baixo consumo e possibilidade de malha local & Garagens com gateway próximo e muitos nós \\
LoRa/LoRaWAN & Longo alcance e mensagens pequenas & Estacionamentos externos ou distribuídos \\
BLE Mesh & Baixo custo e integração com dispositivos próximos & Ambientes menores e com infraestrutura BLE \\
Wi-Fi & Alta disponibilidade e maior taxa de dados & Gateway, câmeras ou locais com alimentação contínua \\
\bottomrule
\end{tabular}
\end{table}

Uma implementação prática pode usar ESP32 ou outro microcontrolador de baixo consumo, um rádio compatível com LoRa ou IEEE 802.15.4 e uma bateria de lítio. O gateway pode ser um computador de placa única ou um equipamento industrial. Para LoRaWAN, o gateway recebe as mensagens e as encaminha ao servidor de rede; para Zigbee, ele atua como coordenador da rede local.

\section{Sensores que podem ser utilizados}
Não existe um único sensor ideal para todos os estacionamentos. O modelo deve considerar preço, precisão, instalação, manutenção, interferência e condições ambientais.

\subsection{Sensor magnético}
Um magnetômetro mede alterações no campo magnético causadas pela massa metálica do veículo. Pode ser instalado no piso e tem a vantagem de ser discreto, consumir pouca energia e funcionar mesmo em ambientes com pouca iluminação. É especialmente interessante para vagas externas e sensores alimentados por bateria. Como limitações, exige calibração e pode sofrer interferência de estruturas metálicas, cabos ou veículos muito próximos. Redes de sensores magnéticos colaborativas já foram estudadas para melhorar a decisão por meio da combinação de leituras de sensores adjacentes \cite{wang2020}.

\subsection{Sensor ultrassônico}
O sensor ultrassônico emite uma onda sonora e mede o tempo de retorno. Instalado acima da vaga, calcula a distância até o veículo; uma distância menor que o limiar indica ocupação. É intuitivo e pode funcionar bem em garagens cobertas. Entretanto, precisa ser protegido contra chuva, poeira e variações de temperatura quando usado ao ar livre, além de exigir posicionamento adequado.

\subsection{Sensor infravermelho}
O sensor infravermelho pode detectar a reflexão da radiação ou a interrupção de um feixe. É barato e simples para protótipos. Estudos comparativos indicam que sensores IR podem apresentar desempenho melhor que sensores baseados apenas em luz em diferentes condições de detecção \cite{alshammari2016}. A luz solar intensa, superfícies muito escuras e obstruções podem afetar a leitura.

\subsection{Sensor de pressão ou presença}
Sensores piezoelétricos, células de carga ou tapetes de pressão detectam a força exercida pelo veículo. São úteis quando o sensor pode ser incorporado ao piso, mas a instalação costuma ser mais invasiva. Também precisam lidar com desgaste mecânico, água e veículos de massas diferentes.

\subsection{Câmera e RFID como recursos complementares}
Câmeras podem confirmar visualmente a ocupação, ler placas ou identificar situações especiais. Porém, demandam mais energia, processamento e cuidados com privacidade. Por isso, uma alternativa eficiente é usar o sensor simples em cada vaga e deixar a câmera apenas em pontos de verificação. RFID não detecta sozinho a presença de qualquer veículo, mas pode identificar veículos autorizados em entradas, saídas ou vagas reservadas. Uma arquitetura com sensores de vaga, leitores RFID e gateway foi discutida por Bagula, Castelli e Zennaro \cite{bagula2015}.

\begin{table}[ht]
\centering
\caption{Comparação simplificada dos sensores.}
\begin{tabular}{p{2.5cm}p{3.2cm}p{3.2cm}p{2.3cm}}
\toprule
Sensor & Vantagem & Limitação & Local indicado \\
\midrule
Magnético & Discreto e econômico & Calibração/interferência & Piso externo \\
Ultrassônico & Medição direta de distância & Instalação e ambiente & Teto de garagem \\
Infravermelho & Baixo custo & Sensível à iluminação & Protótipo/ambiente controlado \\
Pressão & Detecta diretamente o peso & Desgaste e obra no piso & Vaga dedicada \\
Câmera & Vários tipos de análise & Privacidade e energia & Verificação pontual \\
\bottomrule
\end{tabular}
\end{table}

\section{Implementação e fluxo de dados}
O desenvolvimento pode ser dividido em etapas. Primeiro, realiza-se um levantamento das vagas, obstáculos, pontos de energia e área de cobertura. Depois, escolhe-se o sensor e instala-se um protótipo em poucas vagas. Nessa fase devem ser medidos falsos positivos, falsos negativos, alcance do rádio e duração estimada da bateria.

Após a validação, cada nó recebe um identificador único e uma configuração de rede. O gateway registra os dados e os encaminha para uma API. O servidor mantém a última situação de cada vaga, o horário da atualização e o nível de bateria. O aplicativo consulta a API e apresenta as vagas em cores: verde para livre, vermelho para ocupada e cinza para indisponível ou sem atualização recente.

O sistema deve definir um tempo de expiração. Se uma vaga não envia dados por um intervalo, por exemplo cinco minutos, o aplicativo não deve apresentá-la como livre com certeza; deve marcá-la como indisponível ou mostrar que o dado está desatualizado. Essa decisão evita que uma falha de comunicação seja interpretada como uma vaga disponível.

\subsection{Energia, segurança e privacidade}
O consumo pode ser reduzido com sono profundo, leituras periódicas e transmissão apenas quando houver mudança. Em locais com iluminação adequada, painel solar pode complementar a bateria. O gateway deve possuir alimentação mais estável e uma estratégia de armazenamento temporário caso a conexão com a Internet caia.

As mensagens devem ser autenticadas para impedir que um dispositivo falso altere a ocupação das vagas. A comunicação entre aplicativo e servidor deve usar HTTPS, e as chaves dos nós devem ser protegidas. Caso câmeras ou placas sejam utilizadas, o projeto deve limitar a coleta ao necessário, aplicar controle de acesso e definir prazo de retenção das imagens. Para o objetivo básico de informar vagas livres, sensores magnéticos ou ultrassônicos evitam coletar dados pessoais.

\section{Benefícios, limitações e avaliação}
Os benefícios esperados são menor tempo de procura, melhor aproveitamento da infraestrutura, possibilidade de reservar ou priorizar vagas e geração de dados para planejamento. A solução também pode reduzir circulação interna e emissões, embora o efeito dependa da adesão dos usuários e da qualidade da informação.

As principais limitações são o custo inicial, a necessidade de manutenção, a interferência de estruturas metálicas, a cobertura de rádio e a possibilidade de leituras incorretas. A RSSF não elimina a necessidade de sinalização física: o motorista ainda precisa identificar a vaga e respeitar regras de acessibilidade e segurança.

Uma avaliação objetiva deve medir: (i) acurácia da classificação livre/ocupada; (ii) tempo entre a mudança real e a atualização no aplicativo; (iii) taxa de perda de pacotes; (iv) consumo de energia e vida útil estimada; (v) disponibilidade do serviço; e (vi) custo por vaga. O artigo de Chen et al. mostra a importância de avaliar sensores, implantação e distribuição da informação aos usuários em um estacionamento real \cite{chen2016}.

\section{Conclusão}
O estacionamento inteligente é uma aplicação natural para RSSF porque transforma um conjunto grande e distribuído de eventos simples em informação útil para motoristas e administradores. A arquitetura proposta combina nós sensores, comunicação de baixo consumo, gateway, servidor e aplicativo. Sensores magnéticos são uma boa opção para vagas externas e discretas; ultrassônicos são adequados para estruturas cobertas; IR serve bem a protótipos; e câmeras ou RFID podem complementar o sistema quando houver requisitos de verificação ou identificação.

Uma implementação incremental, começando por um pequeno piloto, permite comparar sensores e tecnologias de comunicação antes de instalar toda a rede. O resultado esperado é uma solução escalável, com baixo consumo e informação atualizada, desde que sejam tratados o planejamento de cobertura, a confiabilidade, a segurança e a privacidade.

\section*{Referências}
\addcontentsline{toc}{section}{Referências}
\bibliographystyle{IEEEtran}
\bibliography{referencias}

\end{document}
